\ifdefined\gkver\else\def\gkver{student}\fi
\documentclass[\gkver]{gaokaozhenti}
\begin{document}

\chapter{1991年上海}

\begin{enumerate}
\item
%% number: 1
%% paperName: 1991年上海高考第 1 题 4 分
%% typeId: 1
%% type: 单选题
%% chapter: 原子和原子核
%% point: 半衰期
%% method: 无
%% score: 4
%% degree: 850
%% duplicateId: 0
%% body:
氡 222 衰变为钋 218 的半衰期为 $3.8\Ud$。$20\Ug$ 氡 222 经 $7.6\Ud$ 后还剩下\xzanswer{D}
\fourchoices[answer=D]
{$10\Ug$}
{$5\Ug$}
{$2.6\Ug$}
{$1.25\Ug$}
%% answer: D
%% memo:
\memoanswer{
本题原卷及材料中未提供详解，待补充。
}

\item
%% number: 2
%% paperName: 1991年上海高考第 2 题 4 分
%% typeId: 1
%% type: 单选题
%% chapter: 功和能
%% point: 机械能守恒定律
%% method: 无
%% score: 4
%% degree: 650
%% duplicateId: 0
%% body:
a、b、c 三球自同一高度以相同速率抛出，a 球竖直上抛，b 球水平抛出，c 球竖直下抛。设三球落地时的速率分别为 $v_{a}$、$v_{b}$、$v_{c}$。则\xzanswer{D}
\fourchoices[answer=D]
{$v_{a} > v_{b} > v_{c}$}
{$v_{a} = v_{b} > v_{c}$}
{$v_{a} > v_{b} = v_{c}$}
{$v_{a} = v_{b} = v_{c}$}
%% answer: D
%% memo:
\memoanswer{
本题原卷及材料中未提供详解，待补充。
}

\item
%% number: 3
%% paperName: 1991年上海高考第 3 题 4 分
%% typeId: 1
%% type: 单选题
%% chapter: 气体性质
%% point: 理想气体状态方程
%% method: 无
%% score: 4
%% degree: 650
%% duplicateId: 0
%% body:
一定质量的理想气体处于某一初始状态，现要使它的温度经过状态变化后，回到初始状态的温度，用下列哪个过程可以实现？\xzanswer{A}
\fourchoices[answer=A]
{先保持压强不变而使体积膨胀，接着保持体积不变而减小压强}
{先保持压强不变而使体积减小，接着保持体积不变而减小压强}
{先保持体积不变而增大压强，接着保持压强不变而使体积膨胀}
{先保持体积不变而减小压强，接着保持压强不变而使体积减小}
%% answer: A
%% memo:
\memoanswer{
本题原卷及材料中未提供详解，待补充。
}

\item
%% number: 4
%% paperName: 1991年上海高考第 4 题 4 分
%% typeId: 1
%% type: 单选题
%% chapter: 运动和力
%% point: 动量守恒定律
%% method: 无
%% score: 4
%% degree: 650
%% duplicateId: 0
%% body:
质量相同的三个小球 a、b、c 在光滑水平面上以相同的速率运动，它们分别与原来静止的三个小球 A、B、C 相碰（a 与 A 相碰，b 与 B 相碰，c 与 C 相碰）。碰后，a 球继续沿原来方向运动；b 球静止不动；c 球被弹回而向反方向运动。这时，A、B、C 三球中动量最大的是\xzanswer{C}
\fourchoices[answer=C]
{A 球}
{B 球}
{C 球}
{D 球}
%% answer: C
%% memo:
\memoanswer{
本题原卷及材料中未提供详解，待补充。
}

\item
%% number: 5
%% paperName: 1991年上海高考第 5 题 4 分
%% typeId: 1
%% type: 单选题
%% chapter: 电场
%% point: 带电粒子的平衡
%% method: 无
%% score: 4
%% degree: 450
%% duplicateId: 0
%% body:
如图所示三个点电荷 $q_{1}$、$q_{2}$、$q_{3}$ 固定在一直线上，$q_{2}$ 与 $q_{3}$ 的距离为 $q_{1}$ 与 $q_{2}$ 距离的 2 倍，每个电荷所受的静电力的合力均为零。由此可以判定，三个电荷的电量之比 $q_{1}:q_{2}:q_{3}$ 为\xzanswer{A}
\onepicture[w=9cm]{\includesvg{figs/05}}
\fourchoices[answer=A]
{$-9:4:-36$}
{$9:4:36$}
{$-3:2:-6$}
{$3:2:6$}
%% answer: A
%% memo:
\memoanswer{
本题原卷及材料中未提供详解，待补充。
}

\item
%% number: 6
%% paperName: 1991年上海高考第 6 题 4 分
%% typeId: 1
%% type: 单选题
%% chapter: 磁场
%% point: 安培力
%% method: 无
%% score: 4
%% degree: 650
%% duplicateId: 0
%% body:
如图所示，一根有质量的金属棒 MN，两端用细软导线连接后悬挂于 a、b 两点。棒的中部处于方向垂直纸面向里的匀强磁场中，棒中通有电流，方向从 M 流向 N，此时悬线上有拉力。为了使拉力等于零，可\xzanswer{C}
\onepicture[w=6cm]{\includesvg{figs/06}}
\fourchoices[answer=C]
{适当减小磁感应强度}
{使磁场反向}
{适当增大电流强度}
{使电流反向}
%% answer: C
%% memo:
\memoanswer{
本题原卷及材料中未提供详解，待补充。
}

\item
%% number: 7
%% paperName: 1991年上海高考第 7 题 16 分
%% typeId: 1
%% type: 单选题
%% chapter: 电磁感应
%% point: 切割磁感线电动势
%% method: 无
%% score: 16
%% degree: 450
%% duplicateId: 0
%% body:
图中 PQRS 为一正方形导线框，它以恒定速度向右进入以 MN 为边界的匀强磁场，磁场方向垂直线框平面，MN 线与线框的边成 $45^{\circ}$ 角。E、F 分别为 PS 与 PQ 的中点。关于线框中的感生电流，正确的说法是\xzanswer{B}
\onepicture[w=8cm]{\includesvg{figs/07}}
\fourchoices[answer=B]
{当 E 点经过边界 MN 时，线框中感应电流最大}
{当 P 点经过边界 MN 时，线框中感应电流最大}
{当 F 点经过边界 MN 时，线框中感应电流最大}
{当 Q 点经过边界 MN 时，线框中感应电流最大}
%% answer: B
%% memo:
\memoanswer{
本题原卷及材料中未提供详解，待补充。
}

\item
%% number: 8
%% paperName: 1991年上海高考第 8 题 4 分
%% typeId: 1
%% type: 单选题
%% chapter: 机械振动
%% point: 弹簧振子
%% method: 无
%% score: 4
%% degree: 450
%% duplicateId: 0
%% body:
一平台沿竖直方向作简谐振动，一物体置于振动平台上随台一起运动。当振动平台处于什么位置时，物体对台面的正压力最大？\xzanswer{C}
\fourchoices[answer=C]
{当振动平台运动到最高点时}
{当振动平台向下运动过振动中心点时}
{当振动平台运动到最低点时}
{当振动平台向上运动过振动中心点时}
%% answer: C
%% memo:
\memoanswer{
本题原卷及材料中未提供详解，待补充。
}

\item
%% number: 9
%% paperName: 1991年上海高考第 9 题 5 分
%% typeId: 2
%% type: 多选题
%% chapter: 磁场
%% point: 带电粒子在磁场中的运动
%% method: 无
%% score: 5
%% degree: 650
%% duplicateId: 0
%% body:
在同一匀强磁场中，质子和电子各自在垂直于磁场的平面内作半径相同的匀速圆周运动。质子的质量为 $m_{p}$，电子的质量为 $m_{e}$。则\xzanswer{AC}
\fourchoices[answer=AC]
{质子与电子的速率之比等于 $m_{e}/m_{p}$}
{质子与电子的动量大小之比等于 $m_{e}/m_{p}$}
{质子与电子的动能之比等于 $m_{e}/m_{p}$}
{质子与电子的圆周运动周期之比等于 $m_{e}/m_{p}$}
%% answer: AC
%% memo:
\memoanswer{
本题原卷及材料中未提供详解，待补充。
}

\item
%% number: 10
%% paperName: 1991年上海高考第 10 题 5 分
%% typeId: 2
%% type: 多选题
%% chapter: 曲线运动
%% point: 曲线运动条件
%% method: 无
%% score: 5
%% degree: 650
%% duplicateId: 0
%% body:
如图所示，物体在恒力 $F$ 作用下沿曲线从 A 运动到 B，这时，突然使它所受力的反向，大小不变，即由 $F$ 变为 $-F$。在此力作用下，物体以后的运动情况，下列说法正确的是\xzanswer{ABD}
\onepicture[w=6cm]{\includesvg{figs/10}}
\fourchoices[answer=ABD]
{物体不可能沿曲线 Ba 运动}
{物体不可能沿直线 Bb 运动}
{物体不可能沿曲线 Bc 运动}
{物体不可能沿原曲线由 B 返回 A}
%% answer: ABD
%% memo:
\memoanswer{
本题原卷及材料中未提供详解，待补充。
}

\item
%% number: 11
%% paperName: 1991年上海高考第 11 题 5 分
%% typeId: 2
%% type: 多选题
%% chapter: 气体性质
%% point: 热力学第一定律
%% method: 无
%% score: 5
%% degree: 650
%% duplicateId: 0
%% body:
关于物体内能，下列说法正确的是\xzanswer{BC}
\fourchoices[answer=BC]
{相同质量的两种物体，升高相同的温度，内能增量一定相同}
{一定量 $0\UdC$ 的水结成 $0\UdC$ 的冰，内能一定减少}
{一定量气体体积增大，但既不吸热也不放热，内能一定减少}
{一定量气体吸收热量而保持体积不变，内能一定减少}
%% answer: BC
%% memo:
\memoanswer{
本题原卷及材料中未提供详解，待补充。
}

\item
%% number: 12
%% paperName: 1991年上海高考第 12 题 5 分
%% typeId: 2
%% type: 多选题
%% chapter: 电磁感应
%% point: 电磁感应中图象
%% method: 无
%% score: 5
%% degree: 450
%% duplicateId: 0
%% body:
图（b）中 A、B 为两个相同的环形线圈，共轴并靠近放置。A 线圈中通有如图（a）所示的交流电 $i$，则\xzanswer{ABC}
\onepicture[w=10cm]{\includesvg{figs/12}}
\fourchoices[answer=ABC]
{在 $t_{1}$ 到 $t_{2}$ 时间内 A、B 两线圈相吸}
{在 $t_{2}$ 到 $t_{3}$ 时间内 A、B 两线圈相斥}
{$t_{1}$ 时刻两线圈间作用力为零}
{$t_{2}$ 时刻两线圈间吸力最大}
%% answer: ABC
%% memo:
\memoanswer{
本题原卷及材料中未提供详解，待补充。
}

\item
%% number: 13
%% paperName: 1991年上海高考第 13 题 5 分
%% typeId: 2
%% type: 多选题
%% chapter: 运动和力
%% point: 牛顿定律的应用
%% method: 图象法
%% score: 5
%% degree: 450
%% duplicateId: 0
%% body:
在光滑的水平面上，放着两块长度相同，质量分别为 $M_{1}$ 和 $M_{2}$ 的木板，在两木板的左端各放一个大小、形状、质量完全相同的物块，如图所示。开始时，各物均静止。今在两物块上各作用一水平恒力 $F_{1}$、$F_{2}$，当物块与木板分离时，两木板的速度分别为 $v_{1}$ 和 $v_{2}$。物块与两木板之间的摩擦系数相同。下列说法正确的是\xzanswer{BD}
\onepicture[w=11cm]{\includesvg{figs/13}}
\fourchoices[answer=BD]
{若 $F_{1} = F_{2}$，$M_{1} > M_{2}$，则 $v_{1} > v_{2}$}
{若 $F_{1} = F_{2}$，$M_{1} < M_{2}$，则 $v_{1} > v_{2}$}
{若 $F_{1} > F_{2}$，$M_{1} = M_{2}$，则 $v_{1} > v_{2}$}
{若 $F_{1} < F_{2}$，$M_{1} = M_{2}$，则 $v_{1} > v_{2}$}
%% answer: BD
%% memo:
\memoanswer{
由物块和木板的受力分析判断，物块和木板做匀加速运动，因物块和木板的加速度不同，导致物体与木板分离，利用物块和木板的位移之差求得木板的速度，然后加以讨论。对物块：$F - \mu mg = ma_{1}$，对木板：$\mu mg = Ma_{2}$。从刚开始运动到木块分离有

$\frac{1}{2}a_{1}t^{2} - \frac{1}{2}a_{2}t^{2} = L$

由以上各式得 $t = \sqrt{\frac{2L}{\frac{F}{m} - \mu g - \frac{\mu mg}{M}}}$

讨论：若 $F_{1} = F_{2}$，$M_{1} > M_{2}$，两木板的加速度 $a_{2} < a_{2}'$，运动时间 $t_{1} < t_{2}$，由 $v = at$ 可知 $v_{1} < v_{2}$。选项 A 错误。

若 $F_{1} = F_{2}$，$M_{1} < M_{2}$，两木板的加速度 $a_{2} > a_{2}'$，运动时间 $t_{1} > t_{2}$，由 $v = at$ 可知 $v_{1} > v_{2}$。选项 B 正确。

若 $F_{1} > F_{2}$，$M_{1} = M_{2}$，两木板的加速度 $a_{2} = a_{2}'$，运动时间 $t_{1} < t_{2}$，由 $v = at$ 可知 $v_{1} < v_{2}$。选项 C 错误。

若 $F_{1} < F_{2}$，$M_{1} = M_{2}$，两木板的加速度 $a_{2} = a_{2}'$，运动时间 $t_{1} > t_{2}$，由 $v = at$ 可知 $v_{1} > v_{2}$。选项 D 正确。

正确选项为 BD。
}

\item
%% number: 14
%% paperName: 1991年上海高考第 14 题 4 分
%% typeId: 3
%% type: 填空题
%% chapter: 其他
%% point: 物理学史
%% method: 无
%% score: 4
%% degree: 850
%% duplicateId: 0
%% body:
卢瑟福在 \tkanswer{$\alpha$ 粒子散射} 实验的基础上，提出了原子核式结构学说。\tkanswer{麦克斯韦}在总结前人研究电磁现象的基础上，建立了完整的电磁场理论，预言了电磁波的存在。
%% answer: $\alpha$ 粒子散射，麦克斯韦
\jdanswer{
$\alpha$ 粒子散射，麦克斯韦
}
%% memo:
\memoanswer{
本题原卷及材料中未提供详解，待补充。
}

\item
%% number: 15
%% paperName: 1991年上海高考第 15 题 4 分
%% typeId: 3
%% type: 填空题
%% chapter: 气体性质
%% point: 阿伏伽德罗常数
%% method: 无
%% score: 4
%% degree: 650
%% duplicateId: 0
%% body:
一热水瓶水的质量约 $2.2\Ukg$，它所包含的水分子数目约为 \tkanswer{$7.3\times10^{25}$}。（取二位有效数字。阿伏伽德罗常数取 $6.0\times10^{23}\Umol^{-1}$）
%% answer: $7.3\times10^{25}$
\jdanswer{
$7.3\times10^{25}$
}
%% memo:
\memoanswer{
本题原卷及材料中未提供详解，待补充。
}

\item
%% number: 16
%% paperName: 1991年上海高考第 16 题 4 分
%% typeId: 3
%% type: 填空题
%% chapter: 原子和原子核
%% point: 人工核反应
%% method: 无
%% score: 4
%% degree: 650
%% duplicateId: 0
%% body:
一个锂核（$\ce{^{7}_{3}Li}$）受到一个质子的轰击，变成两个 $\alpha$ 粒子，这一过程的核反应方程是 \tkanswer{$\ce{^{7}_{3}Li + ^{1}_{1}H -> 2 ^{4}_{2}He}$}。已知一个氢原子的质量是 $1.6736\times10^{-27}\Ukg$，一个锂原子的质量是 $11.6505\times10^{-27}\Ukg$，一个氦原子的质量是 $6.6466\times10^{-27}\Ukg$。上述核反应所释放的能量等于 \tkanswer{$2.78\times10^{-12}\UJ$}。
%% answer: $\ce{^{7}_{3}Li + ^{1}_{1}H -> 2 ^{4}_{2}He}$，$2.78\times10^{-12}\UJ$
\jdanswer{
$\ce{^{7}_{3}Li + ^{1}_{1}H -> 2 ^{4}_{2}He}$，$2.78\times10^{-12}\UJ$
}
%% memo:
\memoanswer{
本题原卷及材料中未提供详解，待补充。
}

\item
%% number: 17
%% paperName: 1991年上海高考第 17 题 4 分
%% typeId: 3
%% type: 填空题
%% chapter: 力物体的平衡
%% point: 力矩平衡
%% method: 无
%% score: 4
%% degree: 650
%% duplicateId: 0
%% body:
如图所示一根长为 $L$ 的轻杆 OA，可绕水平轴 O 在竖直平面内自由转动。左端 A 挂一质量为 $m$ 的物体，从杆上一点 B 系一不会伸长的细绳。将绳跨过光滑的钉子 C 与弹簧 K 连接，弹簧处于拉伸状态。已知 $OB = OC = \frac{2}{3}L$，弹簧伸长量恰等于 BC。由此可知，弹簧的倔强系数（劲度系数）等于 \tkanswer{$\frac{9mg}{4L}$}。
\onepicture[align=right, w=7cm]{\includesvg{figs/17}}
%% answer: $\frac{9mg}{4L}$
\jdanswer{
$\frac{9mg}{4L}$
}
%% memo:
\memoanswer{
本题原卷及材料中未提供详解，待补充。
}

\item
%% number: 18
%% paperName: 1991年上海高考第 18 题 4 分
%% typeId: 3
%% type: 填空题
%% chapter: 磁场
%% point: 带电粒子在磁场中的运动
%% method: 无
%% score: 4
%% degree: 450
%% duplicateId: 0
%% body:
如图所示，质量为 $m$，带电量为 $+q$ 的粒子，从两平行电极板正中央垂直电力线和磁力线以速度 $v$ 飞入。已知两板间距为 $d$，磁感应强度为 $B$，这时粒子恰能直线穿过电场和磁场区域（重力不计）。今将磁感应强度增大到某值，则粒子将落到极板上，当粒子落到极板上时的动能为 \tkanswer{$\frac{1}{2}(mv^{2} - qdvB)$}。
\onepicture[align=right, w=7cm]{\includesvg{figs/18}}
%% answer: $\frac{1}{2}(mv^{2} - qdvB)$
\jdanswer{
$\frac{1}{2}(mv^{2} - qdvB)$
}
%% memo:
\memoanswer{
本题原卷及材料中未提供详解，待补充。
}

\item
%% number: 19
%% paperName: 1991年上海高考第 19 题 4 分
%% typeId: 7
%% type: 作图题
%% chapter: 光学
%% point: 透镜
%% method: 无
%% score: 4
%% degree: 650
%% duplicateId: 0
%% body:
将焦距为 $f$ 的凸透镜切成上下两半，沿主轴拉开距离 $f$，如图所示。点光源 S 置于透镜主轴上离左边半个透镜 $f$ 处。该装置可演示两束光的干涉现象。请画出点光源 S 经上、下两半透镜后的光束，并用斜影线画出两束光发生干涉的区域。
\drawpicanswer{\onepicture[align=right, w=10cm]{\includesvg{figs/19}}}{\onepicture[align=right, w=11cm]{\includesvg{figs/19_an}}}
%% answer: 如图所示
\jdanswer{
作图如图所示。
}
%% memo:
\memoanswer{
本题原卷及材料中未提供详解，待补充。
}

\item
%% number: 20
%% paperName: 1991年上海高考第 20 题 4 分
%% typeId: 3
%% type: 填空题
%% chapter: 气体性质
%% point: 气体图象
%% method: 无
%% score: 4
%% degree: 650
%% duplicateId: 0
%% body:
如图表示 $0.2\Umol$ 某种气体的压强与温度关系，图中 $p_{0}$ 为标准大气压。气体在 D 状态时的体积是 \tkanswer{$5.6\UL$}。
\onepicture[align=right, w=7cm]{\includesvg{figs/20}}
%% answer: $5.6\UL$
\jdanswer{
$5.6\UL$
}
%% memo:
\memoanswer{
本题原卷及材料中未提供详解，待补充。
}

\item
%% number: 21
%% paperName: 1991年上海高考第 21 题 4 分
%% typeId: 3
%% type: 填空题
%% chapter: 电路
%% point: 闭合电路欧姆定律
%% method: 无
%% score: 4
%% degree: 450
%% duplicateId: 0
%% body:
图中，AB 和 A$'$B$'$ 是长度均为 $2\Ukm$，每千米电阻值为 $1\UO$ 的两根输电线。今发现在距离 A 和 A$'$ 等远的两点 C 和 C$'$ 间发生漏电，相当于在两点间连接了一个电阻。现用电动势为 $90\UV$，内阻不计的电源接在 AA$'$ 间时，测得 BB$'$ 间电压为 $72\UV$。把此电源接在 BB$'$ 时，测得 AA$'$ 间电压为 $45\UV$，由此可知 A 与 C 相距 \tkanswer{$0.4\Ukm$}。
\onepicture[align=right, w=7.5cm]{\includesvg{figs/21}}
%% answer: $0.4\Ukm$
\jdanswer{
$0.4\Ukm$
}
%% memo:
\memoanswer{
本题原卷及材料中未提供详解，待补充。
}

\item
%% number: 22
%% paperName: 1991年上海高考第 22 题 4 分
%% typeId: 4
%% type: 实验题
%% chapter: 光学
%% point: 光电效应
%% method: 无
%% score: 4
%% degree: 650
%% duplicateId: 0
%% body:
在演示光电效应的实验中，把某种金属板连在验电器上。第一次，用弧光灯直接照射金属板，验电器的指针就张开一个角度。第二次，在弧光灯和金属板之间，插入一块普通的玻璃板，再用弧光灯照射，验电器指针不张开。由此可以判定，使金属板产生光电效应的是弧光中的\xzanswer{B}
\fourchoices[answer=B]
{可见光成份}
{紫外光成份}
{红外光成份}
{无线电波成份}
%% answer: B
\jdanswer{
B
}
%% memo:
\memoanswer{
用弧光灯直接照射金属板，金属板逸出光电子后带正电，使验电器的指针张开。插入普通玻璃板后，因为玻璃板能吸收紫外线，而可见光依然能通过玻璃板照到金属板上，验电器指针不张开，所以，可以判定，使金属板产生光电效应的是弧光中的紫外线成分。
}

\item
%% number: 23
%% paperName: 1991年上海高考第 23 题 5 分
%% typeId: 4
%% type: 实验题
%% chapter: 电路
%% point: 测电动势与内阻
%% method: 无
%% score: 5
%% degree: 650
%% duplicateId: 0
%% body:
用半偏法测图中的电流表 G 的内阻 $R_{g}$。将下列说法中正确的选出来。\xzanswer{ABC}
\onepicture[w=7cm]{\includesvg{figs/23}}
\fourchoices[answer=ABC]
{电键 $K_{1}$ 接通前，$R_{1}$ 必须调节到高阻值处}
{电键 $K_{2}$ 接通前，$R_{2}$ 的阻值无需调节}
{当电流表示数从满偏电流 $I_{g}$ 调到半偏电流 $I_{g}/2$ 时，$R_{2}$ 中电流强度稍大于 $I_{g}/2$}
{$R_{1}$ 阻值的大小，对实验误差没有影响}
%% answer: ABC
\jdanswer{
ABC
}
%% memo:
\memoanswer{
本题原卷及材料中未提供详解，待补充。
}

\item
%% number: 24
%% paperName: 1991年上海高考第 24 题 3 分
%% typeId: 4
%% type: 实验题
%% chapter: 直线运动
%% point: 打点计时器公式
%% method: 无
%% score: 3
%% degree: 650
%% duplicateId: 0
%% body:
用接在 $50\UHz$ 交流低压电源上的打点计时器，测定小车作匀加速直线运动的加速度。某次实验中得到的一条纸带如图所示，从比较清晰的点起，每五个打印点取一个点作为计数点，分别标明 0、1、2、3、4。量得 0 与 1 两点间距离 $S_{1} = 30\Umm$，3 与 4 两点间距离 $S_{4} = 48\Umm$，则小车在 0 与 1 两点间的平均速度为 \tkanswer{$0.30\Ums$}，小车的加速度为 \tkanswer{$0.60\Umsq$}。
\onepicture[align=right, w=11cm]{\includesvg{figs/24}}
%% answer: $0.30\Ums$，$0.60\Umsq$
\jdanswer{
$0.30\Ums$，$0.60\Umsq$
}
%% memo:
\memoanswer{
本题原卷及材料中未提供详解，待补充。
}

\item
%% number: 25
%% paperName: 1991年上海高考第 25 题 6 分
%% typeId: 4
%% type: 实验题
%% chapter: 气体性质
%% point: 查理定律
%% method: 无
%% score: 6
%% degree: 650
%% duplicateId: 0
%% body:
用如图所示的实验装置，研究体积不变时气体的压强与温度的关系，当时大气压为 $H\UcmHg$。封有一定质量气体的烧瓶，浸在冰水混合物中，U 形压强计可动管 A 和固定管 B 中的水银面刚好相平。将烧瓶浸入温度为 $t\UdC$ 的热水中时，B 管水银面将 \tkanswer{下降}，这时应将 A 管 \tkanswer{上升}，（以上二格填“上升”或“下降”），使 B 管中水银面 \tkanswer{回到原处}。记下此时 A、B 两管中水银面的高度差为 $h\Ucm$，此状态下瓶中气体的压强为 \tkanswer{$(H + h)\UcmHg$}。
\onepicture[align=right, w=6.5cm]{\includesvg{figs/25}}
%% answer: 下降，上升，回到原处，$(H + h)\UcmHg$
\jdanswer{
下降，上升，回到原处，$(H + h)\UcmHg$
}
%% memo:
\memoanswer{
本题原卷及材料中未提供详解，待补充。
}

\item
%% number: 26
%% paperName: 1991年上海高考第 26 题 7 分
%% typeId: 4
%% type: 实验题
%% chapter: 电路
%% point: 分压与分流
%% method: 无
%% score: 7
%% degree: 650
%% duplicateId: 0
%% body:
一电灯，电阻为 $4R$。设计一个电路，使电灯两端的电压调节范围尽可能大。可利用的器材为：一个电动势为 $E$；内阻为 $R/5$ 的电源；一个最大阻值为 $R$ 的滑线变阻器；一个电键；导线若干。

（a）将设计的电路图画在下方虚线框内。

（b）灯两端电压可调节的范围为 \tkanswer{$0\sim0.8E$}。

（c）按设计的电路图，将如图所示仪器用线（表示导线）连接起来。
\onepicture[align=right, w=7cm]{\includesvg{figs/26}}
%% answer: （a）画图如图。（b）$0\sim0.8E$。（c）见图。
\jdanswer{
（a）设计的电路如图所示。

\onepicture[w=6cm]{\includesvg{figs/26_an1}}

（b）$0\sim0.8E$。

（c）连线如图所示。

\onepicture[w=8cm]{\includesvg{figs/26_an2}}
}
%% memo:
\memoanswer{
本题原卷及材料中未提供详解，待补充。
}

\item
%% number: 27
%% paperName: 1991年上海高考第 27 题 12 分
%% typeId: 6
%% type: 计算题
%% chapter: 光学
%% point: 光的折射
%% method: 无
%% score: 12
%% degree: 650
%% duplicateId: 0
%% body:
某三棱镜的横截面是一直角三角形，如图所示，$\angle A = 90^{\circ}$，$\angle B = 30^{\circ}$，$\angle C = 60^{\circ}$；棱镜材料的折射率为 $n$，底面 BC 涂黑。入射光沿平行于底面 BC 的方向射向 AB 面，经 AB 面和 AC 面折射后出射。
\begin{enumerate}
	\item
	求出射光线与入射光线延长线间的夹角 $\delta$。
	\item
	为使上述入射光线能从 AB 面出射，折射率 $n$ 的最大值为多少？
\end{enumerate}
\onepicture[align=right, w=9cm]{\includesvg{figs/27}}
%% answer: （1）$\delta = \sin^{-1}\frac{\sqrt{4n^{2} - 3}}{2} - 30^{\circ}$；（2）$n \le \frac{\sqrt{7}}{2}$
\jdanswer{
\begin{enumerate}
	\item
	$\delta = \sin^{-1}\frac{\sqrt{4n^{2} - 3}}{2} - 30^{\circ}$

	\item
	$n \le \frac{\sqrt{7}}{2}$

\end{enumerate}
}
%% memo:
\memoanswer{
\onepicture[align=right, w=9cm]{\includesvg{figs/27_an}}

（1）如图，设光在 AB 面上入射角为 $i$，折射角为 $\alpha$，在 AC 面上入射角为 $\beta$，出射角为 $\gamma$，由折射定律

$\sin i = n\sin\alpha$

$i = 60^{\circ}$，$\sin\alpha = \frac{1}{n}\sin60^{\circ} = \frac{\sqrt{3}}{2n}$　　2 分

$\alpha + \beta = 90^{\circ}$，$\sin\beta = \cos\alpha$

$\sin\beta = \sqrt{1 - \sin^{2}\alpha} = \frac{\sqrt{4n^{2} - 3}}{2n}$　　2 分

由折射定律

$n\sin\beta = \sin\gamma$

$\sin\gamma = n\cdot\frac{\sqrt{4n^{2} - 3}}{2n} = \frac{\sqrt{4n^{2} - 3}}{2}$　　2 分

$\delta = \sin^{-1}\frac{\sqrt{4n^{2} - 3}}{2} - 30^{\circ}$　　2 分

（2）要使有光线从 AC 面射出，应有

$\sin\gamma \le 1$

即 $\frac{\sqrt{4n^{2} - 3}}{2} \le 1$　　2 分

得 $n \le \frac{\sqrt{7}}{2}$　　2 分
}

\item
%% number: 28
%% paperName: 1991年上海高考第 28 题 10 分
%% typeId: 6
%% type: 计算题
%% chapter: 电路
%% point: 电功电功率
%% method: 无
%% score: 10
%% degree: 450
%% duplicateId: 0
%% body:
某电炉在额定电压下的功率为 $P_{0} = 300\UW$。某电源在不接负载时的路端电压与电炉的额定电压相同。当将电炉接到该电源上时，电炉实际消耗的功率为 $P_{1} = 324\UW$。若将两个这样的电炉并联接入该电源，两个电炉实际消耗的总功率 $P_{2}$ 为多少？
%% answer: $P_{2} = 535.5\UW$
\jdanswer{
$P_{2} = 535.5\UW$
}
%% memo:
\memoanswer{
设电炉额定电压为 $U$，电阻为 $R$，电源内阻为 $r$，电源电动势即为 $U$。故

$P_{0} = \frac{U^{2}}{R} = 400$ ①，2 分

$P_{1} = \frac{U^{2}}{(R + r)^{2}}R = 324$ ②，2 分

$P_{2} = \frac{U^{2}}{(\frac{R}{2} + r)^{2}}\cdot\frac{R}{2}$ ③，2 分

由①②消去 $U$，解得

$\frac{r}{R} = \sqrt{\frac{P_{0}}{P_{1}}} - 1 = \sqrt{\frac{400}{324}} - 1 = \frac{1}{9}$

代入③式

$P_{2} = \frac{U^{2}}{R}\frac{1}{2(\frac{1}{2} + \frac{r}{R})^{2}} = \frac{400}{2(\frac{1}{2} + \frac{1}{9})^{2}} = 535.5\UW$　　4 分
}

\item
%% number: 29
%% paperName: 1991年上海高考第 29 题 14 分
%% typeId: 6
%% type: 计算题
%% chapter: 曲线运动
%% point: 平抛运动
%% method: 无
%% score: 14
%% degree: 450
%% duplicateId: 0
%% body:
如图所示，长为 $l$ 的轻绳一端系于固定点 O，另一端系质量为 $m$ 的小球。将小球从 O 点正下方 $l/4$ 处，以一定初速水平向右抛出，经一定时间绳被拉直，以后小球将以 O 为支点在竖直平面内摆动。已知绳刚被拉直时，绳与竖直线成 $60^{\circ}$ 角。求：
\begin{enumerate}
	\item
	小球水平抛出时的初速 $v_{0}$。
	\item
	在绳被拉紧的瞬间，支点 O 受到的冲量 $I$。
	\item
	小球摆到最低点时，绳所受的拉力 $T$。
\end{enumerate}
\onepicture[align=right, w=8cm]{\includesvg{figs/29}}
%% answer: （1）$v_{0} = \frac{\sqrt{6gl}}{2}$；（2）$I = m\sqrt{2gl}$；（3）$T = 2mg$
\jdanswer{
\begin{enumerate}
	\item
	$v_{0} = \frac{\sqrt{6gl}}{2}$

	\item
	$I = m\sqrt{2gl}$

	\item
	$T = 2mg$

\end{enumerate}
}
%% memo:
\memoanswer{
\onepicture[align=right, w=8cm]{\includesvg{figs/29_an}}

（1）小球在绳被拉直前作平抛运动，如图所示。设小球抛出后经时间 $t$ 绳被拉直

$v_{0}t = l\sin60^{\circ}$　　1 分

$\frac{1}{2}gt^{2} = l\cos60^{\circ} - \frac{l}{4}$　　1 分

由此解得

$t = \sqrt{\frac{l}{2g}}$

$v_{0} = \frac{\sqrt{6gl}}{2}$　　2 分

（2）在绳被拉直前瞬时，小球速度的水平分量为 $v_{0}$，竖直分量为 $gt$，速度大小

$v = \sqrt{v_{0}^{2} + g^{2}t^{2}}$

以 $v_{0} = \frac{\sqrt{6gl}}{2}$，$gt = g\sqrt{\frac{l}{2g}} = \frac{\sqrt{2gl}}{2}$ 代入

得 $v = \sqrt{2gl}$　　2 分

速度与竖直方向的夹角 $\varphi = \tan^{-1}\frac{v_{0}}{gt} = \tan^{-1}\sqrt{3} = 60^{\circ}$　　3 分

可见小球速度与绳沿同一线，小球动量在绳拉力的冲量作用下减速为零，于是 O 点所受冲量即为绳拉直前一瞬时小球的动量：

$I = Ft = mv = m\sqrt{2gl}$　　2 分

（3）以后小球作摆动，由机械能守恒，小球到最低点时动能

$\frac{1}{2}mv^{2} = mgl - mgl\cos60^{\circ} = \frac{1}{2}mgl$　　1 分

由牛顿定律 $T - mg = \frac{mv^{2}}{l}$　　1 分

得绳拉力 $T = mg + \frac{mv^{2}}{l} = 2mg$　　1 分
}

\end{enumerate}

\end{document}
